\section{مباحث نظری}
\label{sec:theory}

\subsection{پارامترهای تنظیم \lr{PostgreSQL}}

در کنار مستندات رسمی، دو منبع دیگر هم برای تحقیق درباره‌ی این پارامترها مفیدند: \lr{postgresqlco.nf} \cite{pgconfig} که توضیح و مقدار پیش‌فرض هر پارامتر را کنار هم نشان می‌دهد، و \lr{PgTune} \cite{pgtune} که با گرفتن مشخصات سخت‌افزار و نوع بار کاری، یک نقطه‌ی شروع برای این مقادیر پیشنهاد می‌دهد.

\paragraph{\lr{shared\_buffers}.} این پارامتر اندازه‌ی کش مشترک پیج‌های پایگاه‌داده را در حافظه‌ی خود \lr{PostgreSQL} تعیین می‌کند \cite{pgresource}. بزرگ‌کردن زیاد آن همیشه بهتر نیست: چون سیستم‌عامل هم پیج‌ها را در کش خودش نگه می‌دارد، عددی بین \lr{25\%} تا \lr{40\%} حافظه‌ی کل معمولاً خوب است. بیشتر از این، همان داده دوبار (یک‌بار در کش \lr{PostgreSQL}، یک‌بار در کش سیستم‌عامل) نگه داشته می‌شود بدون فایده‌ی اضافه.

\paragraph{\lr{work\_mem}.} برخلاف \lr{shared\_buffers}، این مقدار برای هر عملیات \lr{Sort} یا \lr{Hash} در هر کوئری جداگانه است، نه یک‌بار برای کل سیستم. یک کوئری پیچیده با چند \lr{Sort}/\lr{Hash} می‌تواند چند برابر این مقدار حافظه مصرف کند. پس مقداری که برای یک اتصال خوب است، ممکن است برای صدها اتصال هم‌زمان کل حافظه‌ی سرور را تمام کند. در این پروژه، مقدار سراسری پایین نگه داشته شده و فقط در سطح نشست برای کوئری‌های سنگین بالا می‌رود.

\paragraph{\lr{effective\_cache\_size}.} این پارامتر هیچ حافظه‌ای تخصیص نمی‌دهد؛ فقط یک تخمین به برنامه‌ریز می‌دهد از این‌که چقدر داده احتمالاً در کش جا می‌شود. عدد بزرگ‌تر باعث می‌شود برنامه‌ریز هزینه‌ی \lr{I/O} تصادفی را کمتر تخمین بزند و بیشتر به سمت ایندکس‌ها برود.

\paragraph{\lr{wal\_compression} و \lr{default\_toast\_compression}.} هر دو یک مبادله بین \lr{CPU} و \lr{I/O} دیسک هستند. \lr{lz4} سریع‌تر از الگوریتم پیش‌فرض \lr{pglz} است و روی حجم کاری بارگذاری انبوه این پروژه، هزینه‌ی \lr{CPU} اضافی آن خیلی کمتر از فایده‌ی کاهش نوشتار دیسک است.

\paragraph{\lr{autovacuum}.} \lr{PostgreSQL} از کنترل هم‌روندی چندنسخه‌ای (\lr{MVCC}) استفاده می‌کند: هر \lr{UPDATE}/\lr{DELETE} نسخه‌ی قدیمی سطر را پاک نمی‌کند، فقط آن را «مرده» علامت می‌زند تا تراکنش‌های دیگر که هنوز به آن نیاز دارند خراب نشوند. \lr{autovacuum} این سطرهای مرده را پاک می‌کند. آستانه‌ی پیش‌فرض آن (درصدی از اندازه‌ی جدول) برای یک صف کار که خیلی سریع‌تر از رشد جدول، سطر مرده تولید می‌کند مناسب نیست؛ به همین دلیل جدول \lr{message\_queue} با آستانه‌ی مطلق (تعداد ثابت سطر) تنظیم شده، نه آستانه‌ی نسبی.

\paragraph{\lr{maintenance\_work\_mem}.} این پارامتر هم مثل \lr{work\_mem} یک محدودیت حافظه است، ولی نه برای کوئری‌های معمولی؛ برای عملیات‌های نگه‌داری مثل \lr{CREATE INDEX}، \lr{VACUUM} و \lr{ALTER TABLE} استفاده می‌شود. چون این عملیات‌ها معمولاً یکی‌یکی و کم‌تعداد اجرا می‌شوند، نه هم‌زمان برای صدها اتصال مثل \lr{work\_mem}، می‌شود مقدار آن را خیلی بزرگ‌تر گذاشت بدون نگرانی از تمام‌شدن حافظه‌ی سرور. مقدار پیش‌فرض \lr{PostgreSQL} فقط \lr{64MB} است؛ برای پروژه‌ای که باید ایندکس یک جدول \lr{90} میلیون سطری بسازد، این عدد کم است -- عدد بزرگ‌تر یعنی مرتب‌سازی ساخت ایندکس در حافظه انجام شود، نه با چند پاس روی دیسک. به همین دلیل در این پروژه \lr{768MB} گذاشته شده.

\paragraph{\lr{wal\_buffers}.} قبل از این‌که رکوردهای \lr{WAL} (لاگ پیش‌نوشت، برای بازیابی بعد از کرش) واقعاً روی دیسک نوشته شوند، اول در یک بافر مشترک در حافظه جمع می‌شوند؛ این پارامتر اندازه‌ی همان بافر است. مقدار پیش‌فرض آن \lr{-1} است، یعنی خودکار تنظیم شود (معمولاً حدود \lr{1/32} از \lr{shared\_buffers}). بزرگ‌ترکردن آن به تراکنش‌های هم‌زمان زیاد کمک می‌کند، چون رکوردهای بیشتری قبل از یک نوشتار واقعی جمع می‌شوند؛ ولی بعد از یک حد مشخص (اندازه‌ی چیزی که بین دو چک‌پوینت واقعاً نوشته می‌شود) بزرگ‌ترکردن بیشتر فایده‌ای ندارد. در این پروژه روی \lr{16MB} تنظیم شده.

\paragraph{\lr{min\_wal\_size} و \lr{max\_wal\_size}.} این دو، فضایی که \lr{PostgreSQL} برای فایل‌های \lr{WAL} نگه می‌دارد را کنترل می‌کنند و مستقیماً روی زمان‌بندی چک‌پوینت اثر می‌گذارند. رسیدن حجم \lr{WAL} به \lr{max\_wal\_size} باعث می‌شود یک چک‌پوینت زودتر از موعد اجرا شود. \lr{max\_wal\_size} کوچک یعنی چک‌پوینت‌های مکرر -- هزینه‌ی \lr{I/O} بیشتر ولی بازیابی سریع‌تر بعد از کرش؛ \lr{max\_wal\_size} بزرگ یعنی چک‌پوینت‌های کم‌تر -- برای بارگذاری انبوه بهتر است، ولی بازیابی بعد از کرش کندتر می‌شود. مقدار پیش‌فرض \lr{PostgreSQL 16} برای این دو به‌ترتیب \lr{80MB} و \lr{1GB} است. در این پروژه، عمداً یک \lr{max\_wal\_size} نسبتاً کوچک (\lr{4GB}) انتخاب شده تا خط پایه رفتار واقعی یک بارگذاری \lr{I/O}-محدود را نشان دهد.

\paragraph{\lr{seq\_page\_cost} و \lr{random\_page\_cost}.} این دو، واحد هزینه‌ای هستند که برنامه‌ریز کوئری برای مقایسه‌ی پلن‌های مختلف استفاده می‌کند: \lr{seq\_page\_cost} هزینه‌ی خواندن یک صفحه به‌صورت پی‌درپی است (پیش‌فرض \lr{1.0})، و \lr{random\_page\_cost} هزینه‌ی خواندن یک صفحه به‌صورت تصادفی (پیش‌فرض \lr{4.0}). این نسبت \lr{4} به \lr{1} از دنیای دیسک‌های مکانیکی می‌آید، جایی که یک جهش تصادفی سر دیسک واقعاً چند برابر کندتر از خواندن پی‌درپی بود. روی \lr{SSD}/\lr{NVMe} این تفاوت خیلی کمتر است، پس نگه‌داشتن \lr{random\_page\_cost} روی \lr{4} باعث می‌شود برنامه‌ریز به‌اشتباه فکر کند ایندکس همیشه گران‌تر از اسکن کامل است. در این پروژه \lr{random\_page\_cost} روی \lr{1.1} تنظیم شده تا با واقعیت سخت‌افزار سازگار باشد.

\paragraph{\lr{enable\_seqscan}، \lr{enable\_indexscan} و \lr{enable\_bitmapscan}.} این سه، کلید قطع‌ووصل سطح-نشست هستند؛ خاموش‌کردن هرکدام آن روش دسترسی را کاملاً ممنوع نمی‌کند، فقط هزینه‌ی آن را مصنوعی خیلی بالا می‌برد تا برنامه‌ریز تقریباً همیشه از آن فرار کند. کاربرد اصلی‌شان تست و مقایسه است، نه تنظیم دائمی در \lr{postgresql.conf}: مثلاً می‌شود موقتاً \lr{enable\_bitmapscan} را خاموش کرد تا ببینیم اگر ایندکس در دسترس نبود پلن چه شکلی می‌شد، بدون این‌که واقعاً ایندکس را حذف کنیم.

ما همین کار را واقعاً برای سناریوی \lr{1} (بخش \ref{sec:indexing}) انجام دادیم: با \lr{enable\_bitmapscan = off} و \lr{enable\_indexscan = off}، برنامه‌ریز مجبور شد به‌جای \lr{Bitmap Heap Scan} (که خودش انتخاب کرده بود، در \lr{8953.3} میلی‌ثانیه) به \lr{Parallel Seq Scan} برگردد -- این‌بار در \lr{14522.6} میلی‌ثانیه، یعنی حتی کندتر. نکته‌ی مهم: این عدد را نباید مستقیم با عدد اصلی «قبل از ایندکس» (\lr{4381.2} میلی‌ثانیه) مقایسه کرد، چون آن اندازه‌گیری ساعت‌ها زودتر و با کش تازه‌تری از دیسک انجام شده بود؛ اما مقایسه‌ی این دو عدد در همین نشست (هردو با وضعیت کش یکسان) نشان می‌دهد که برنامه‌ریز، با توجه به کش فعلی، واقعاً بهترین گزینه‌ی موجود را انتخاب کرده بود -- فقط آن گزینه از عدد خط پایه‌ی اصلی بدتر بود.

\subsection{مقایسه‌ی انواع ایندکس}

\begin{itemize}
\item \textbf{\lr{B-Tree}} (پیش‌فرض \lr{PostgreSQL}): یک ساختار درختی که کلیدها را مرتب نگه می‌دارد. برای تساوی، بازه (\lr{<}, \lr{>}, \lr{BETWEEN}) و مرتب‌سازی خوب است. فقط ستون‌های اول یک ایندکس چندستونی قابل‌استفاده‌اند.
\item \textbf{\lr{Hash}}: فقط برای تساوی (\lr{=}) خوب است؛ نه بازه و نه مرتب‌سازی را پشتیبانی نمی‌کند. از \lr{PostgreSQL 10} به بعد بعضی وقت‌ها کوچک‌تر و سریع‌تر از \lr{B-Tree} است، اما این مزیت به‌ندرت به‌اندازه‌ای است که ارزش از دست‌دادن بازه/مرتب‌سازی را داشته باشد -- به همین دلیل در این پروژه استفاده نشده.
\item \textbf{\lr{GiST}} (\lr{Generalized Search Tree}): برخلاف بقیه که برای یک نوع داده‌ی خاص طراحی شده‌اند، \lr{GiST} یک چارچوب عمومی است. با نوشتن چندتابع خاص (تشخیص هم‌پوشانی، فشرده‌سازی، ادغام)، می‌شود آن را برای داده‌های خیلی متفاوت به‌کار برد: داده‌ی هندسی/جغرافیایی (مثلاً در \lr{PostGIS})، بازه‌ها (\lr{range types})، جست‌وجوی نزدیک‌ترین همسایه (\lr{KNN}، با عملگر \lr{<->})، و حتی جست‌وجوی متن به‌عنوان جایگزین \lr{GIN}. ساخت و به‌روزرسانی آن معمولاً سریع‌تر از \lr{GIN} است، ولی چون بعضی از پیاده‌سازی‌هایش \lr{lossy} هستند (ممکن است نتیجه‌ی اضافی هم برگردد که باید دوباره چک شود)، جست‌وجو در آن کمی کندتر و اندازه‌اش بزرگ‌تر از یک \lr{GIN} هم‌ارز است.
\item \textbf{\lr{GIN}} (\lr{Generalized Inverted Index}): برای مقادیری که هر سطر چندتا از آن‌ها دارد -- آرایه، \lr{JSONB}، جست‌وجوی متن. برای هر مقدار داخلی، یک لیست از سطرهایی که آن مقدار را دارند نگه می‌دارد. برای «آیا این مقدار توی این ستون هست؟» (\lr{@>}, \lr{\&\&}) خیلی سریع است، اما مفهوم ترتیب ندارد و فقط با \lr{Bitmap Index Scan} در دسترس است، نه \lr{Index Scan} ساده. با پسوند \lr{btree\_gin}، یک \lr{GIN} می‌تواند چند ستون از نوع‌های مختلف را با هم ترکیب کند و -- برخلاف \lr{B-Tree} -- با فیلتر روی هرکدام از این ستون‌ها به‌تنهایی هم کار کند.
\item \textbf{\lr{BRIN}} (\lr{Block Range Index}): برای ستون‌هایی که ترتیب فیزیکی‌شان روی دیسک با ترتیب مقدارشان همخوانی دارد (مثلاً یک ستون زمانی که به‌ترتیب درج شده) خیلی کوچک و کم‌هزینه است، چون فقط بازه‌ی مقدار هر بلوک دیسک را نگه می‌دارد، نه هر سطر را. در این پروژه استفاده نشده چون داده‌ی \lr{IMDb} بعد از عبور از صف پیام، دیگر لزوماً با همان ترتیبی که \lr{BRIN} لازم دارد درج نمی‌شود.
\end{itemize}

\paragraph{کدام ایندکس برای \lr{JSONB} مناسب است؟} از این چهار نوع، \lr{GIN} انتخاب استاندارد \lr{PostgreSQL} برای \lr{JSONB} است -- همان‌طور که بالا گفته شد، دقیقاً برای سؤال «آیا این کلید/مقدار توی این \lr{JSONB} هست؟» (\lr{@>}, \lr{?}) طراحی شده. \lr{GiST} هم روی \lr{JSONB} قابل‌استفاده است ولی رایج نیست و معمولاً کندتر از \lr{GIN} عمل می‌کند. \lr{B-Tree} و \lr{Hash} اصلاً مناسب این کار نیستند، چون فقط تساوی/ترتیب کل مقدار \lr{JSONB} را می‌فهمند، نه محتوای داخلش را.

\subsection{\lr{JSONB} و مکانیزم \lr{TOAST}}

\paragraph{تفاوت \lr{JSON} و \lr{JSONB}.} نوع \lr{JSON} متن خام ورودی را دقیقاً همان‌طور که نوشته شده نگه می‌دارد -- با همان فاصله‌ها، همان ترتیب کلیدها، و حتی کلیدهای تکراری. هر بار که این مقدار خوانده یا پرس‌وجو شود، \lr{PostgreSQL} باید کل متن را از نو تجزیه (\lr{parse}) کند. \lr{JSONB} برعکس این است: در زمان نوشتن یک‌بار تجزیه و به یک قالب باینری آماده تبدیل می‌شود. نتیجه یک مبادله‌ی ساده است: درج \lr{JSONB} کمی کندتر است (چون باید ساخته شود)، ولی خواندن و جست‌وجو خیلی سریع‌تر است (چون تجزیه‌ی دوباره لازم نیست) -- برای همین تقریباً همیشه \lr{JSONB} به \lr{JSON} ترجیح داده می‌شود.

\paragraph{\lr{JSONB} چگونه در داخل \lr{PostgreSQL} ذخیره می‌شود؟} مقدار \lr{JSONB} به‌صورت یک قالب باینریِ از قبل تجزیه‌شده و «تجزیه‌شده» (\lr{decomposed binary format}) ذخیره می‌شود، نه به‌صورت متن. این قالب یک هدر به همراه آرایه‌ای از افست‌ها/طول‌ها دارد که مستقیماً به محل هر کلید و مقدار در همان مقدار باینری اشاره می‌کند. کلیدهای هر شیء بر اساس طول و سپس ترتیب بایتی مرتب نگه داشته می‌شوند، پس پیداکردن یک کلید مشخص با جست‌وجوی دودویی (\lr{binary search}) روی همین آرایه ممکن است -- بدون این‌که کل مقدار از اول خوانده شود. سه چیز در این فرایند تبدیل از بین می‌رود: فاصله‌های خالی متن اصلی، ترتیب اصلی کلیدها (چون بر اساس طول/بایت دوباره مرتب می‌شوند)، و کلیدهای تکراری (اگر یک کلید دوبار در ورودی بیاید، فقط آخرین مقدارش نگه داشته می‌شود). اعداد هم به‌جای متن، در یک نوع عددی داخلی ذخیره می‌شوند. همین ساختار باینریِ با افستِ مستقیم است که هم جست‌وجوی سریع را ممکن می‌کند و هم پایه‌ی ساخت ایندکس \lr{GIN} روی \lr{JSONB} است.

\paragraph{عملگرها و توابع مهم \lr{JSONB}.} چندتای پرکاربردش: \lr{->} یک فیلد/عضو را به‌صورت \lr{jsonb} برمی‌گرداند؛ \lr{->>} همان را به‌صورت متن برمی‌گرداند؛ \lr{@>} می‌پرسد «آیا سمت چپ شامل سمت راست است؟» -- همان عملگر \lr{containment} که ایندکس \lr{GIN} برایش ساخته می‌شود؛ \lr{?} می‌پرسد «آیا این کلید در سطح بالای \lr{JSONB} وجود دارد؟». از توابع هم \lr{jsonb\_set()} (تغییر یک مقدار داخلی بدون خواندن-تغییر-نوشتن کل شیء در کد برنامه)، \lr{jsonb\_agg()} (جمع‌کردن چند سطر در یک آرایه‌ی \lr{JSONB}) و \lr{jsonb\_each()} (باز کردن یک \lr{JSONB} به سطرهای کلید-مقدار) پرکاربردند. در این پروژه، ستون \lr{title\_ratings.metadata} دقیقاً با عملگر \lr{@>} قابل‌پرس‌وجوست -- مثلاً \lr{metadata @> '\{"popular": true\}'} که روی داده‌ی واقعی این پروژه \lr{17{,}906} سطر برمی‌گرداند.

\lr{TOAST} (\lr{The Oversized-Attribute Storage Technique}) \cite{pgtoast} مکانیزمی است که به \lr{PostgreSQL} اجازه می‌دهد مقادیر متغیرطول بزرگ (مثل \lr{JSONB}، متن یا آرایه) را بیرون از سطر اصلی، در یک جدول کمکی نگه دارد. چهار روش ذخیره برای هر ستون وجود دارد: \lr{PLAIN} (بدون فشرده‌سازی یا انتقال؛ برای انواع با اندازه‌ی ثابت)، \lr{MAIN} (فشرده‌سازی درون‌خطی، انتقال به بیرون فقط اگر واقعاً لازم باشد)، \lr{EXTENDED} (پیش‌فرض \lr{JSONB}؛ هم فشرده‌سازی هم انتقال به بیرون مجاز است) و \lr{EXTERNAL} (بدون فشرده‌سازی ولی انتقال به بیرون مجاز -- برای وقتی که دسترسی تصادفی به بخشی از مقدار مهم‌تر از فضاست). فعال‌شدن این مکانیزم -- همان‌طور که در بخش \ref{sec:queue} گفتیم -- بر اساس کل سطر است، نه هر ستون؛ همین نکته پایه‌ی طراحی بسته‌بندی دسته‌ای این پروژه است.

\paragraph{ساخت ایندکس \lr{GIN} روی \lr{JSONB}.} دستور ساده است: \lr{CREATE INDEX ... ON table USING GIN (jsonb\_column)}. حالت پیش‌فرض (\lr{jsonb\_ops}) هم کلیدها هم مقدارها را ایندکس می‌کند و از \lr{@>}, \lr{?}, \lr{?|}, \lr{?\&} پشتیبانی می‌کند. یک گزینه‌ی دیگر، \lr{jsonb\_path\_ops}، فقط از \lr{@>} پشتیبانی می‌کند ولی ایندکس کوچک‌تر و سریع‌تری می‌سازد -- برای ستونی مثل \lr{metadata} در این پروژه که (اگر روزی لازم شود) فقط قرار است با \lr{containment} پرس‌وجو شود، \lr{jsonb\_path\_ops} انتخاب بهتری از پیش‌فرض بود.

\paragraph{\lr{JSONB} یا جدول نرمال‌شده؟} قاعده‌ی ساده: اگر قرار است روی یک فیلد داخلی زیاد فیلتر/پیوند/گروه‌بندی کنید، بهتر است آن فیلد یک ستون واقعی با ایندکس \lr{B-Tree} باشد -- برنامه‌ریز آمار بهتری از یک ستون واقعی دارد تا از داخل یک \lr{JSONB}. اما اگر ساختار داده متغیر است (هر سطر ممکن است کلیدهای متفاوتی داشته باشد)، یا قرار است کل شیء با هم خوانده/نوشته شود، \lr{JSONB} بهتر است چون از ساختن چندین جدول کمکی با ستون‌های اغلب \lr{NULL} جلوگیری می‌کند. ستون \lr{metadata} در این پروژه دقیقاً نمونه‌ی دوم است: داده‌ی مشتق‌شده‌ای که هیچ‌کدام از هشت سناریو روی آن فیلتر نمی‌کند، پس نیازی به نرمال‌سازی ندارد.

\subsection{نظریه‌ی صف پیام}

\paragraph{الگوی \lr{Competing Consumers}.} در این الگو، چند مصرف‌کننده‌ی مستقل از یک صف مشترک کار برمی‌دارند تا پردازش سریع‌تر شود. مشکل اصلی این است که یک کار نباید توسط بیش از یک مصرف‌کننده برداشته شود، بدون این‌که مصرف‌کننده‌ها مجبور شوند پشت سر هم صف بکشند.

\paragraph{مکانیزم \lr{LISTEN}/\lr{NOTIFY}.} این دو دستور، مکانیزم داخلی \lr{PostgreSQL} برای اطلاع‌رسانی بین نشست‌ها هستند. یک نشست با \lr{LISTEN channel\_name} روی یک کانال ثبت‌نام می‌کند؛ هر نشست دیگری با \lr{NOTIFY channel\_name, 'payload'} (یا تابع \lr{pg\_notify}) یک پیام کوتاه متنی به همه‌ی نشست‌های ثبت‌نام‌شده روی آن کانال می‌فرستد -- این پیام معمولاً بعد از \lr{commit} همان تراکنش، به‌صورت غیرهم‌زمان به گیرنده‌ها می‌رسد.

\paragraph{چرا \lr{LISTEN}/\lr{NOTIFY} به‌تنهایی برای صف کافی نیست.} مشکل اصلی این است که \lr{NOTIFY} هیچ حافظه‌ای ندارد: اگر لحظه‌ی ارسال هیچ نشستی روی آن کانال گوش نمی‌داد، آن پیام برای همیشه گم می‌شود -- نه صف‌بندی، نه امکان بازپخش. همچنین \lr{NOTIFY} به‌صورت پخش (\lr{broadcast}) به همه‌ی گوش‌دهنده‌ها می‌رسد، نه فقط یکی؛ یعنی خودش هیچ تضمینی نمی‌دهد که یک پیام دقیقاً یک‌بار توسط یکی از مصرف‌کننده‌ها برداشته شود -- دقیقاً همان مشکلی که \lr{Competing Consumers} باید حلش کند. به همین دلیل این پروژه، حالت واقعی صف (وضعیت، مهلت، شمارنده‌ی تلاش) را در جدول \lr{message\_queue} نگه می‌دارد، نه در \lr{NOTIFY}. اگر لازم شود، می‌شود \lr{LISTEN}/\lr{NOTIFY} را فقط به‌عنوان یک زنگ‌خطر سبک اضافه کرد -- مصرف‌کننده به‌جای \lr{polling} دوره‌ای منتظر \lr{NOTIFY} بماند -- ولی وضعیت واقعی هر پیام همچنان باید در جدول بماند.

\paragraph{\lr{SKIP LOCKED} در برابر قفل معمولی.} یک \lr{SELECT ... FOR UPDATE} معمولی، اگر سطرهایش قبلاً قفل شده باشند، منتظر می‌ماند تا آزاد شوند -- یعنی مصرف‌کننده‌ها پشت سر هم صف می‌کشند. اضافه‌کردن \lr{SKIP LOCKED} \cite{pglocking} این رفتار را عوض می‌کند: هر سطر قفل‌شده از نتیجه حذف می‌شود و تراکنش می‌رود سراغ سطر آزاد بعدی. نتیجه این است که \lr{N} مصرف‌کننده‌ی هم‌زمان، \lr{N} دسته‌ی کاملاً جدا از صف برمی‌دارند، بدون این‌که منتظر هم بمانند -- دقیقاً همان چیزی که \lr{Competing Consumers} نیاز دارد.

\paragraph{تضمین‌های تحویل.} سه سطح رایج وجود دارد: \lr{at-most-once} (پیام ممکن است هیچ‌وقت پردازش نشود، ولی هیچ‌وقت هم دوبار پردازش نمی‌شود)، \lr{at-least-once} (پیام حتماً حداقل یک‌بار پردازش می‌شود، ولی ممکن است در بعضی حالت‌ها -- مثلاً خرابی مصرف‌کننده قبل از تأیید -- بیش از یک‌بار هم پردازش شود)، و \lr{exactly-once} (که در سیستم‌های توزیع‌شده‌ی واقعی، بدون یک لایه‌ی هماهنگ‌کننده‌ی اضافه، تقریباً غیرممکن یا خیلی گران است). صف این پروژه تضمین \lr{at-least-once} می‌دهد: مکانیزم مهلت با \lr{locked\_until} تضمین می‌کند اگر مصرف‌کننده‌ای قبل از \lr{ack} از کار بیفتد، پیام (بعد از تمام‌شدن مهلت و اجرای \lr{reap\_expired}) دوباره به صف برمی‌گردد -- ولی این یعنی ممکن است همان پیام دوبار پردازش شود.

\paragraph{بی‌اثری (\lr{Idempotency}).} چون تضمین \lr{at-least-once} اجازه می‌دهد پردازش تکراری اتفاق بیفتد، مصرف‌کننده باید طوری طراحی شود که پردازش دوباره‌ی یک پیام، نتیجه‌ی نهایی را عوض نکند. در این پروژه، این کار با \lr{ON CONFLICT DO NOTHING} روی کلید اصلی هر جدول انجام شده: اگر یک دسته دوبار درج شود (مثلاً به‌خاطر برگشتن به صف)، تلاش دوم سطرهای تکراری را نادیده می‌گیرد، بدون خطا و بدون تغییر داده.

\paragraph{جدول-به‌عنوان-صف در برابر راه‌حل تخصصی.} همان‌طور که پروژه‌ی \lr{Postgres For Everything} \cite{pgeverything} هم نشان می‌دهد، \lr{PostgreSQL} می‌تواند نقش خیلی از ابزارهای تخصصی -- از جمله صف پیام -- را بازی کند. استفاده از یک جدول ساده به‌جای یک \lr{message broker} تخصصی (مثل \lr{RabbitMQ} یا \lr{Kafka}) چند مزیت دارد: هیچ زیرساخت جدیدی لازم نیست، و از همه مهم‌تر، می‌شود \lr{enqueue} کردن یک پیام و تغییر بقیه‌ی داده‌ها را در یک تراکنش واحد انجام داد -- یک \lr{broker} جدا این هماهنگی رایگان را نمی‌دهد. عیبش این است که زیر بار خیلی سنگین (هزاران تولیدکننده/مصرف‌کننده هم‌زمان) به‌اندازه‌ی یک \lr{broker} تخصصی مقیاس‌پذیر نیست، و هر بار \lr{poll} یک کوئری واقعی روی همان پایگاه‌داده‌ای است که به کارهای دیگر هم سرویس می‌دهد؛ ویژگی‌های پیشرفته‌تر (مسیریابی پیچیده، بازپخش از یک نقطه‌ی مشخص، داشبورد نظارتی) هم باید از صفر ساخته شوند. برای حجم کار این پروژه -- یک بارگذاری انبوه، نه پیام‌رسانی بین چند سرویس مستقل -- سادگی جدول-به‌عنوان-صف کاملاً کافی و حتی ترجیح داده می‌شود.
